\section{The contradiction}\label{sec:contradiction}

\begin{idea}[What this section does]
We carry out the plan of \cref{sec:ledger}. The graph $G$ only needs
properties (G1) and (G2) from the end of \cref{sec:reduction}. All the
checks below are short ``two-step path'' arguments in the grid. The one
non-trivial input is the pruning lemma.
\end{idea}

\begin{proposition}[Proposition~4.1 of the original]\label{prop:no-concave}
There is no nonempty oriented graph $G$ in which every vertex has an
in-neighbor and
\begin{equation}\label{eq:concave-again}
 \abs U+\abs{F^2(U)}<2\abs{F(U)}\qquad\text{for every }\varnothing\ne U\subsetneq V(G).
\end{equation}
\end{proposition}

\begin{proof}
Suppose $G$ is such a graph, let $X=V(G)$, and use the arc relation of $G$ in
the definitions of \cref{sec:pairs}. By \cref{lem:diag-compatible} we can fix
an admissible pair $(P,Q)$ of maximum value. Write
\[
 \mu=\abs P+\abs Q,\qquad d=\unc(P,Q)=\abs X^2-\abs{\Ff P}-\abs{\Fs Q}.
\]
The proof is a sequence of claims.

\begin{claim}\label{cl:lines-proper}
Every column $P_j$ and every row $Q_i$ is a nonempty proper subset of $X$.
\end{claim}

\begin{proof}
$j\in P_j$ and $i\in Q_i$ because $\Delta\subseteq P\cap Q$. Suppose $P_j=X$. Every
vertex has an in-neighbor, so $F(X)=X$ and column $j$ of $\Ff P$ is all of
$X$ by~\eqref{eq:columnwise}. Choose $w\to j$. Then $(w,w)\in Q$ gives
$(w,j)\in\Fs Q$, while $(w,j)\in\Ff P$, contradicting compatibility. If $Q_i=X$, then
row $i$ of $\Fs Q$ is all of $X$. For $w\to i$, $(w,w)\in P$ gives $(i,w)\in\Ff P$, again a
contradiction.
\end{proof}

\begin{claim}\label{cl:fiber-sums}
$\abs P+\abs{\Ff^2P}<2\abs{\Ff P}$ and $\abs Q+\abs{\Fs^2Q}<2\abs{\Fs Q}$.
\end{claim}

\begin{proof}
Apply~\eqref{eq:concave-again} to each column $U=P_j$, which is allowed by
\cref{cl:lines-proper}, and sum over the $\abs X\ge1$ columns
using~\eqref{eq:columnwise}:
\begin{equation}\label{eq:fiber-sum}
 \abs P+\abs{\Ff^2P}=\sum_j\bigl(\abs{P_j}+\abs{F^2(P_j)}\bigr)<\sum_j2\abs{F(P_j)}=2\abs{\Ff P}.
\end{equation}
The same argument over the rows of $Q$ gives the second inequality.
\end{proof}

Now define the new families as complements of the second iterated images:
\[
 \Rf=(X\times X)\setminus\Fs^2Q,\qquad\Cf=(X\times X)\setminus\Ff^2P .
\]
Note the swap: the new \emph{left} family is built from the old \emph{right} one.

\begin{claim}[Gain]\label{cl:gain}
$\abs\Rf+\abs\Cf>\mu+2d$.
\end{claim}

\begin{proof}
By \cref{cl:fiber-sums},
$\abs\Rf+\abs\Cf=2\abs X^2-\abs{\Fs^2Q}-\abs{\Ff^2P}
>2\abs X^2+\abs P+\abs Q-2\abs{\Ff P}-2\abs{\Fs Q}=\mu+2d$.
\end{proof}

\begin{figure}[ht]
\centering
\begin{tikzpicture}[scale=1.0]
 \begin{scope}
  \node[font=\small] at (0,0.75) {$j$}; \node[font=\small] at (2.0,0.75) {$s$}; \node[font=\small] at (4.0,0.75) {$t$};
  \node[font=\small] at (-0.8,0) {$u$}; \node[font=\small] at (-0.8,-1.5) {$p$}; \node[font=\small] at (-0.8,-3.0) {$i$};
  \foreach \x/\y in {2.0/0,0/-1.5,2.0/-1.5,4.0/-1.5,2.0/-3.0} \node[gridpt] at (\x,\y) {};
  \node[Lpt,label={[font=\footnotesize]right:$(u,s)\in P$}] (us) at (2.0,0) {};
  \node[Hpt,label={[font=\footnotesize,hcol!80!black]below right:$h$}] (h) at (2.0,-1.5) {};
  \node[Rpt,label={[font=\footnotesize]below:$c\in\Cf$}] (c) at (2.0,-3.0) {};
  \node[Rpt,label={[font=\footnotesize]below:$(p,t)\in Q$}] (pt) at (4.0,-1.5) {};
  \node[Lpt,label={[font=\footnotesize]below:$r\in\Rf$}] (r) at (0,-1.5) {};
  \draw[-{Stealth[length=2mm]},thick,famL] (us) -- node[right,font=\footnotesize,black] {$u{\to}p$} (h);
  \draw[-{Stealth[length=2mm]},thick,famL,dashed] (h) -- node[left,font=\footnotesize,black] {$p{\to}i$} (c);
  \draw[-{Stealth[length=2mm]},thick,famR] (pt) -- node[above,font=\footnotesize,black] {$t{\to}s$} (h);
  \draw[-{Stealth[length=2mm]},thick,famR,dashed] (h) -- node[above,font=\footnotesize,black] {$s{\to}j$} (r);
  \node[font=\small,align=center] at (2.0,-4.3) {(a) why $h\notin\Ff P\cup\Fs Q$};
 \end{scope}
 \begin{scope}[xshift=8.2cm]
  \node[font=\small] at (0,0.7) {$v$}; \node[font=\small] at (2.1,0.7) {$s$};
  \node[font=\small] at (-0.8,0) {$s$}; \node[font=\small] at (-0.8,-1.4) {$v$}; \node[font=\small] at (-0.8,-2.8) {$i$};
  \foreach \x/\y in {2.1/0,2.1/-2.8,0/-1.4} \node[gridpt] at (\x,\y) {};
  \node[Lpt,label={[font=\footnotesize]right:$(s,s)\in P$}] (ss) at (2.1,0) {};
  \node[gridpt,draw=black] (vs) at (2.1,-1.4) {};
  \node[Rpt,label={[font=\footnotesize]below:$(i,s)\in\Cf$}] (is) at (2.1,-2.8) {};
  \node[Lpt,label={[font=\footnotesize]below:$(v,v)\in\Rf$}] (vv) at (0,-1.4) {};
  \draw[-{Stealth[length=2mm]},thick,famL] (ss) -- node[right,font=\footnotesize,black] {$s{\to}v$} (vs);
  \draw[-{Stealth[length=2mm]},thick,famL] (vs) -- node[right,font=\footnotesize,black] {$v{\to}i$} (is);
  \node[font=\small,align=center] at (1.0,-4.3) {(b) diagonal members\\ are conflict-free};
 \end{scope}
\end{tikzpicture}
\caption{The two-step path arguments of \cref{cl:diag-free,cl:H-uncovered},
drawn in the grid (rows $=$ first coordinate). (a) For a conflict
$(r,c)$ with opposite corner $h=(p,s)$: if $h\in\Ff P$ via $(u,s)\in P$, then one
more vertical step puts $c$ in $\Ff^2P$; if $h\in\Fs Q$ via $(p,t)\in Q$, then one
more horizontal step puts $r$ in $\Fs^2Q$. Both are impossible. (b) A
conflict of $(v,v)$ with $(i,s)$ would put $(i,s)$ in $\Ff^2P$ via the
diagonal point $(s,s)$.}
\label{fig:twostep}
\end{figure}

\begin{claim}[The diagonal survives]\label{cl:diag-free}
$\Delta\subseteq\Rf\cap\Cf$, and no diagonal member of $\Rf$ or $\Cf$ is in a conflict.
\end{claim}

\begin{proof}
\emph{$\Delta\subseteq\Rf$.} If $(v,v)\in\Fs^2Q$, there are $s\to v$ with
$(v,s)\in\Fs Q$. But $(s,s)\in P$ and $s\to v$ give $(v,s)\in\Ff P$, contradicting
compatibility. \emph{$\Delta\subseteq\Cf$.} If $(v,v)\in\Ff^2P$, there are $p\to v$ with
$(p,v)\in\Ff P$, while $(p,p)\in Q$ gives $(p,v)\in\Fs Q$. Again a contradiction.

\emph{Diagonal members are conflict-free.} If the left member $(v,v)$ conflicted
with $(i,s)\in\Cf$, then $v\to i$ and $s\to v$, so the path $s\to v\to i$ moves
$(s,s)\in P$ to $(i,s)\in\Ff^2P$ (\cref{fig:twostep}(b)), contradicting
$(i,s)\in\Cf$. If $(p,j)\in\Rf$ conflicted with the right member $(v,v)$, then $p\to v$
and $v\to j$, and the path $p\to v\to j$ moves $(p,p)\in Q$ to $(p,j)\in\Fs^2Q$,
contradicting $(p,j)\in\Rf$.
\end{proof}

Let $Z$ and $H$ be the conflict points and opposite corners of $(\Rf,\Cf)$
(\cref{def:ZH}).

\begin{claim}[Opposite corners were uncovered]\label{cl:H-uncovered}
$H\cap(\Ff P\cup\Fs Q)=\varnothing$. Hence $\abs H\le d$.
\end{claim}

\begin{proof}
Let $(p,j)\in\Rf$ and $(i,s)\in\Cf$ conflict, so $p\to i$ and $s\to j$, with
$h=(p,s)$. If $h\in\Ff P$, pick $(u,s)\in P$ with $u\to p$. The path $u\to p\to i$
puts $(i,s)\in\Ff^2P$, contradicting $(i,s)\in\Cf$. If $h\in\Fs Q$, pick $(p,t)\in Q$
with $t\to s$. The path $t\to s\to j$ puts $(p,j)\in\Fs^2Q$, contradicting
$(p,j)\in\Rf$ (\cref{fig:twostep}(a)). Since $\Ff P$ and $\Fs Q$ are disjoint,
$X^2\setminus(\Ff P\cup\Fs Q)$ has exactly $d$ points, so $\abs H\le d$.
\end{proof}

\emph{Pruning.} Apply \cref{lem:pruning} to $\Rf,\Cf$ and the arc relation of $G$.
We obtain $P'\subseteq\Rf$ and $Q'\subseteq\Cf$ satisfying \ref{P1}--\ref{P3}. Let
$d'=\unc(P',Q')$.

\begin{claim}\label{cl:new-admissible}
$(P',Q')$ is admissible, and
$\bigl\lvert Z\setminus(\Ff P'\cup\Fs Q')\bigr\rvert\le d'$.
\end{claim}

\begin{proof}
$P',Q'$ are compatible by \ref{P1}. They contain $\Delta$ by
\cref{cl:diag-free} and \ref{P2}. By \cref{lem:conflict-geometry}(b),
$\Ff P'\cap\Fs Q'=\varnothing$, so exactly $d'$ points of $X^2$ are outside
$\Ff P'\cup\Fs Q'$, and in particular at most $d'$ points of $Z$ are.
\end{proof}

\emph{Conclusion.} By \ref{P3} and \cref{cl:H-uncovered,cl:new-admissible},
at most $\abs H+d'\le d+d'$ members were deleted. With \cref{cl:gain}:
\[
 \val(P',Q')=\abs{P'}+\abs{Q'}+d'\ \ge\ \abs\Rf+\abs\Cf-(d+d')+d'\ >\ \mu+2d-d=\val(P,Q).
\]
This contradicts the maximality of $\val(P,Q)$ among admissible pairs.
\end{proof}

\begin{proof}[Proof of \cref{thm:main}]
If the theorem failed, \cref{prop:reduction} would give a graph that
\cref{prop:no-concave} rules out.
\end{proof}

\begin{remark}[Where the hypotheses were used]\label{rem:hypotheses}
\emph{Orientation} (no opposite arcs) is used four times: in
$\Ntwo(v)=F^2\{v\}\setminus F\{v\}$; in $C\subseteq\Ntwo(v)$ in
\cref{lem:subset-deficit}; to make the blow-up oriented; and in
\cref{lem:diag-compatible} ($\Delta$ is compatible with itself).
\emph{Positive indegrees} of $G$ are used in \cref{cl:lines-proper} (lines
are proper) and in the case $S=V(D)$ of \cref{lem:blowup}.
\emph{Strictness} of~\eqref{eq:concave-again} is used only in
\cref{cl:gain}. Without it the gain would be $\ge\mu+2d$, and the final
inequality would no longer be strict. The pruning lemma itself uses none
of these hypotheses.
\end{remark}

\begin{remark}[Stage 3 as an algorithm]\label{rem:algorithm}
The proof of \cref{prop:no-concave} can be \emph{run} on any oriented graph
$G$ with positive indegrees, whether or not~\eqref{eq:concave-again} holds.
Start with $(\Delta,\Delta)$. Given an admissible pair, form $\Rf,\Cf$,
build $\Gamma$, compute a maximum matching and the vertex cover of
\cref{cl:konig}, and delete the members in the cover. Everything here is
polynomial in $\abs X$; the ``fewest conflict edges'' condition of
\cref{sec:strategy} is used only to \emph{prove} the matching bound, not
to construct the pruning.
\Cref{cl:lines-proper,cl:diag-free,cl:H-uncovered,cl:new-admissible} do
not use~\eqref{eq:concave-again}, so the new pair is admissible, and the
ledger together with the exact identity of \cref{rem:error-term} gives
\[
 \val(P',Q')\ \ge\ \abs\Rf+\abs\Cf-d\ =\ \val(P,Q)-\sum_j\kappa(P_j)-\sum_i\kappa(Q_i).
\]
So as long as every line is strictly concave, the value strictly increases,
and since $\val\le3\abs X^2$ this can happen at most $3\abs X^2$ times.
When the iteration stops, some column $P_j$ or row $Q_i$ has $\kappa\ge0$:
a nonempty proper set $U$ with $\abs U+\abs{F^2U}\ge2\abs{FU}$, found in
polynomial time. If $G=D[T_m]$ is the blow-up of \cref{lem:blowup}, the
master inequality~\eqref{eq:blowup-master} applied to this $U$ shows that
its support $S$ is a nonempty proper subset of $V(D)$ with
$\abs{F^2S\setminus FS}\ge\abs{FS\setminus S}$ (the case $S=V(D)$ always
has $\kappa<0$, and for $S\ne V(D)$ an integer bracket with
$m\cdot\text{bracket}+n\ge0$ and $m>n$ is nonnegative). Thus Stages~2
and~3 together are a polynomial-time algorithm that, given an oriented
graph with positive indegrees, finds a set $S$ violating the subset
deficit~\eqref{eq:subset-deficit}. We do not know how to turn this into a
polynomial-time algorithm for a Seymour \emph{vertex}: Stage~1 argues by
minimality and would have to be unwound recursively. The animation in
\texttt{animation/} runs exactly this iteration on a $4$-vertex graph.
\end{remark}
